\section{Exact metric coordinates}\label{sec:coordinates}

\subsection{Moving backgrounds and the full deformation space}
We first justify the geometric identifications used in the theorem.
For each fixed stable factor $F_i$, solve the HE equation on its
determinant line with a normalized scalar potential. On trace-free
Hermitian endomorphisms of $F_i$, the HE linearization has no kernel:
a kernel element is holomorphic and hence scalar, and its pointwise
trace is zero. The implicit-function theorem and elliptic regularity
therefore give a smooth family $h_i(p)$ of factor HE metrics.
These metrics make the factors polystable; their unchanged scalar
holomorphic endomorphism algebras exclude a nontrivial polystable
decomposition. They consequently remain stable. This is the
moving-factor construction of \cite[Proposition~5.1]{wuebben15}.
The HE existence and polystability correspondence used here are the
classical ones \cite{donaldson,uy,kobayashi}.

The underlying Dolbeault complex remains fixed as $p$ varies.
Its cohomology dimensions are constant, so its harmonic projections
and Green operators vary smoothly. On a smaller parameter
neighborhood their complementary positive spectra have a common
gap. Project a basis of each central harmonic factor to the moving
harmonic space and normalize by the positive square root of its
Gram matrix. Tensor these identifications with the multiplicity
Hom spaces. This gives the smooth $G$-equivariant maps $\iota_p$
used in \eqref{eq:kuranishi}, preserving the trace constraint.
Every uniform assertion below is at a prescribed finite spatial
and parameter regularity order on a common smaller neighborhood.

Smooth positive bundle isometries from $h_p$ to $h_0$ identify the
Hermitian endomorphism bundles. They commute with multiplicity
automorphisms. The range $\ha$ of $P_p$ is therefore fixed.
Since $P_pP_0=P_0$ and $P_0P_p=P_p$, the map
\begin{equation}\label{eq:complement-trivialization}
 T_p=Q_p|_{Q_0},\qquad T_p^{-1}=Q_0|_{Q_p},
\end{equation}
identifies the complementary spaces. These are smooth,
$K$-equivariant identity-plus-finite-rank maps, bounded in every
H\"older space used in the proof.

The convergent construction \eqref{eq:kuranishi} is defined on an
ambient ball, including nonintegrable points. On its zero locus
$d+a_p(x)$ is integrable. We use the full construction as in
\cite[Theorem~8.7]{lam} and \cite[Lemma~2.1]{wuebben15}.
One consequence needed for large reconstructing gauges deserves
an explicit proof.

\begin{lemma}\label{lem:endpoint}
Whenever $x$ and $gx$ belong to the common reconstruction ball,
\[
 a_p(gx)=g\,a_p(x)g^{-1}.
\]
This holds for every $g\in G$, without a bound on its norm.
In particular, integrability is preserved between these endpoints.
\end{lemma}

\begin{proof}
Parallel multiplicity automorphisms commute with the Hodge operators,
including nonunitary automorphisms. Expand the convergent solution
of \eqref{eq:kuranishi} into homogeneous terms
$a_p(x)=\sum_{j\ge1}a_{p,j}(x)$. Each term is obtained recursively
from the equivariant bilinear operation
$-d_p^*G_{2,p}(\,\cdot\,\wedge\,\cdot\,)$, and is therefore
$G$ equivariant for every $g$. At the two endpoints both series
converge. The fixed linear operator $\Ad_g$ commutes with their
sum, proving the identity. The same argument applies to
$\kappa_p$; alternatively the identity conjugates the two
Dolbeault operators. No intermediate point on a group path is
required to lie in the small ball.
\end{proof}

Every sufficiently small integrable deformation has such a harmonic
presentation. Impose $d_p^*a'=0$ by a small determinant-one complex
gauge orthogonal to the full parallel algebra. The linearized
Coulomb equation on that complement is $-d_p^*d$, with a uniformly
bounded inverse. Its right side already lies in the complement.
The implicit-function theorem gives a common local slice.
Integrability and Hodge decomposition then give
$a'=x-d_p^*G_{2,p}(a'\wedge a')$, so contraction uniqueness identifies
$a'$ with $a_p(x)$. This construction uses a fixed elliptic
complement, not the varying tangent space to a finite orbit.

\subsection{The exact sliced equation}
For small $u\in Q_p h^{19/8}(\Herm_0E_0)$ put $h(u)=h_pe^u$.
Even when $d_x=d+a_p(x)$ is not integrable, it determines a unique
compatible unitary connection with this $(0,1)$ part.
In the holomorphic presentation its $(1,0)$ part is
\begin{equation}\label{eq:chern}
 D'_{0,p}+e^{-u}D'_{0,p}(e^u)
                    -e^{-u}a_p(x)^{*h_p}e^u.
\end{equation}
Define the centered curvature
\begin{equation}\label{eq:centered}
 \M(p,x,u)=e^{u/2}
       \bigl(i\Lambda_pF_{d_x,h(u)}-\kappa_E(p)I\bigr)e^{-u/2}.
\end{equation}
It is Hermitian and pointwise trace free, because the determinant
connection remains the fixed HE determinant. At $x=u=0$,
$\M=-b_p$.

For clarity, write $\mathsf C_u,\mathsf S_u,\mathsf B_u$ for
functions of the pointwise commutator $\ad_u$, reserving $C_x$
for the infinitesimal group action:
\begin{equation}\label{eq:exponential-functions}
 \mathsf C_u=\cosh(\ad_u/2),\qquad
 \mathsf S_u=\frac{\sinh(\ad_u/2)}{\ad_u/2},\qquad
 \mathsf B_u=\mathsf S_u^{-1}\mathsf C_u,\qquad
 R_p(u)=P_p\mathsf B_u|_{\ha}.
\end{equation}
The removable singularities have value one. In a unitary basis
diagonalizing $u$, the multipliers of $\mathsf S_u$ and
$\mathsf B_u$ on an off-diagonal entry are $\sinh(a)/a$ and
$a\coth(a)$, where $a=(u_i-u_j)/2$. These are positive and
$a\coth(a)\ge1$. Thus $\mathsf B_u\ge I$ and
\begin{equation}\label{eq:R-inverse}
 R_p(u)\ge I,\qquad \|R_p(u)^{-1}\|_{L^2(\ha)\to L^2(\ha)}\le1.
\end{equation}
This inverse is unrelated to the spectrum of $C_x^*C_x$.

\begin{proposition}\label{prop:slice}
Define
\begin{equation}\label{eq:slice}
 \begin{split}
 \beta(p,x,u)&=-R_p(u)^{-1}P_p\mathsf S_u^{-1}\M(p,x,u),\\
 \dot x&=C_x\beta(p,x,u),\\
 \dot u&=-2\mathsf S_u^{-1}\M(p,x,u)-2\mathsf B_u\beta(p,x,u),\\
 \dot g&=\beta(p,x,u)g.
 \end{split}
\end{equation}
For classical solutions, these equations preserve $P_pu=0$,
trace and determinant. If $x(0)=g(0)x_b$ and $x_b\in Z_p$,
then $x(s)=g(s)x_b$ and $x(s)\in Z_p$ as long as the
coefficients stay in the reconstruction domain. The metric
$H(s)=g(s)^*(h_pe^{u(s)})$ on $d_{x_b}$ satisfies exactly
\begin{equation}\label{eq:hym}
 H^{-1}\partial_sH=-2(i\Lambda_pF_H-\kappa_E(p)I).
\end{equation}
\end{proposition}

\begin{proof}
Projection of $\dot u$ onto $\ha$ is zero by the definition of
$R_p$ and $\beta$. The coefficient $\beta$ is Hermitian and has
weighted trace zero, so the last equation remains in $G$.
Uniqueness for the linear representation equation gives
$x(s)=g(s)x_b$. Lemma~\ref{lem:endpoint} gives
$d_{x(s)}=g(s)d_{x_b}g(s)^{-1}$ and exact integrability.

Transport the metric velocity by $g^{-1}$ and center by
$e^{-u/2}$. Since $\dot gg^{-1}=\beta=\beta^{*h_p}$, the two
gauge derivatives contribute $2\mathsf C_u\beta$. The exponential
derivative contributes
\[
 e^{-u/2}(d\exp)_u(\dot u)e^{-u/2}
   =\int_{-1/2}^{1/2}e^{t\ad_u}\dot u\,dt
   =\mathsf S_u\dot u.
\]
Their sum is $-2\M$ by \eqref{eq:slice}. Curvature transforms by
the same bundle isometry, proving \eqref{eq:hym}.
The calculation uses neither commutativity of $u$ and $\beta$
nor commutativity of successive values of $\beta$.
\end{proof}

The classical calculation requires one time derivative and two
spatial derivatives of $u$. The parabolic realization below
supplies this regularity for positive times and continuity at
the initial time in $h^{17/8}$.
For constant $k\in K$, the transformation
$(x,u,g)\mapsto(kx,kuk^{-1},kg)$ preserves \eqref{eq:slice}
and reconstructs the same metric. This is the unitary redundancy
of these coordinates. Removing $g$ instead would identify
distinct parallel positive metrics at $x=0$.

\subsection{Relative metrics and large gauges}
A right-invariant Riemannian metric on the finite-dimensional
group $G$ is complete. Its distance satisfies
\[
 d_R(g,g_*)=d_R(gg_*^{-1},I).
\]
The order of the two factors matters. This right quotient,
rather than $g_*^{-1}g$, is what cancels the absolute gauge
size in the comparison of metrics.

\begin{lemma}\label{lem:relative}
Let $\Psi:B_r(V)\to Q_ph^\tau(\Herm_0E_0)$ be a $C^1$ map,
where $\tau>0$ is nonintegral, with
$\|\Psi(x)\|_{h^\tau}\le Cr^2$ and
$\|D\Psi(x)\|_{V\to h^\tau}\le Cr$.
Consider metrics on the same $d_{x_b}$,
\[
 H=g^*(h_pe^{\Psi(x)+v}),\qquad
 H_*=g_*^*(h_pe^{\Psi(x_*)}),\qquad x=gx_b,\quad x_*=g_*x_b,
\]
with $x,x_*\in B_r$ and $\|v\|_{h^\tau}$ and $d_R(g,g_*)$
sufficiently small. The centered logarithm
\[
 \log\!\left(e^{-\Psi(x_*)/2}q^{*h_p}
             e^{\Psi(x)+v}q e^{-\Psi(x_*)/2}\right),
 \qquad q=gg_*^{-1},
\]
has $h^\tau$ norm at most
$C(\|v\|_{h^\tau}+d_R(g,g_*))$. In particular,
$e^{-\eta}H_*\le H\le e^\eta H_*$ for this same bound $\eta$.
The constants are independent of the absolute gauge norms
and of the stabilizers.
\end{lemma}

\begin{proof}
In an identity chart, $\|q-I\|\le C d_R(g,g_*)$.
Since $x=qx_*$, it follows that
$|x-x_*|\le Cr\,d_R(g,g_*)$, without inverting $C_x$.
The segment between $x$ and $x_*$ lies in the ambient ball,
so the graph derivative bound gives
$\|\Psi(x)-\Psi(x_*)\|_{h^\tau}\le Cr^2d_R(g,g_*)$.
The matrix $q$ is parallel for the split connection, with
uniformly bounded spatial norms in the identity chart.
Multiplication, exponential and logarithm are smooth operations
in the little-H\"older algebra. The displayed logarithm vanishes
at $q=I,v=0$, which gives its norm estimate. Its $C^0$
operator bound implies the metric-order statement after
pullback by $g_*$. No estimate on $g$ or $g_*$ enters.
\end{proof}
